%% agentic-coding-loop.tex: the agentic coding cycle (Chapter 5).
%% House conceptual style. The palette echoes the evolutionary loop of Chapter 10
%% (violet execute, green read, orange commit) so the two cycles read as kin. The
%% verification oracle hangs off Execute as a run--fail--fix feedback arc (the
%% grounded loop that lets Execute converge); Read is annotated as the human
%% ownership checkpoint, the one step that cannot be delegated.
%% _concept-style.tex: shared TikZ styles for the course's CONCEPTUAL block diagrams.
%% The leading underscore marks this as the one file in figures/ that is not a
%% figure; every other figures/*.tex holds exactly one figure body.
%%
%% \input this at the top of every conceptual figure (before \begin{tikzpicture})
%% so each such diagram speaks one visual language: rounded "block" nodes, stealth
%% arrows, and natural `<hue>!15' fills. Redefining these styles on each \input is
%% idempotent and harmless. The reference for the house parameters is
%% resources/STYLE-figures.md. Figure~\ref{figure:AgentLoop} is the exemplar.
%%
%% Scope: use for conceptual/relational diagrams (boxes-and-arrows). Do NOT use for
%% product mock-ups (the rig figures have their own shared language) or for
%% data/plot figures.
\tikzset{
  % the standard block node: geometry only — apply a `fill=<hue>!15' per node.
  conceptbox/.style={draw, rounded corners=4pt, minimum width=2.4cm,
                     minimum height=0.85cm, align=center, font=\small},
  % a flow arrow between blocks.
  conceptflow/.style={->, >=stealth', thick},
  % an edge/annotation label.
  conceptlabel/.style={font=\scriptsize, align=center},
}%

\begin{tikzpicture}
% Four-node cycle, clockwise from top-left; the oracle is a satellite of Execute.
\node[conceptbox, fill=blue!15]   (specify) at (0, 2)   {Specify};
\node[conceptbox, fill=violet!15] (exec)    at (5, 2)   {Execute};
\node[conceptbox, fill=green!15]  (read)    at (5, 0)   {Read};
\node[conceptbox, fill=orange!15] (commit)  at (0, 0)   {Commit\\or revise};
\node[conceptbox, fill=gray!15, minimum width=3.0cm] (oracle) at (10.8, 2)
    {Verification oracle\\\scriptsize interpreter, tests};

\draw[conceptflow] (specify) -- node[conceptlabel, midway, above]{instruction} (exec);
\draw[conceptflow] (exec) -- node[conceptlabel, midway, right]{diff} (read);
\draw[conceptflow] (read) -- node[conceptlabel, midway, below]{understood} (commit);
\draw[conceptflow] (commit) -- node[conceptlabel, midway, left]{revise / next change} (specify);
\draw[conceptflow, <->] (exec) -- node[conceptlabel, midway, above]{run / fail / fix} (oracle);

\node[conceptlabel] at (5, -0.75) {\emph{ownership checkpoint}};
\end{tikzpicture}%
